\documentclass[11pt,a4paper]{article}

\usepackage[margin=2.3cm]{geometry}
\usepackage[portuguese]{babel}
\usepackage{booktabs}
\usepackage{xcolor}
\usepackage{amsmath}
\usepackage{enumitem}
\usepackage{parskip}
\usepackage{titlesec}
\usepackage{listings}

\definecolor{azul}{HTML}{1E3A8A}
\definecolor{cinza}{HTML}{F4F6F8}
\definecolor{linha}{HTML}{CBD5E1}

\titleformat{\section}{\large\bfseries\color{azul}}{\thesection.}{0.5em}{}
\titlespacing*{\section}{0pt}{14pt}{6pt}

\lstset{
  basicstyle=\ttfamily\small,
  backgroundcolor=\color{cinza},
  frame=single,
  rulecolor=\color{linha},
  framexleftmargin=4pt,
  xleftmargin=8pt,
  aboveskip=6pt,
  belowskip=6pt,
  columns=fullflexible,
  keepspaces=true
}

\begin{document}

\begin{center}
  {\LARGE\bfseries\color{azul} Assignment-1 --- MAXSAT / SAT}\\[3pt]
  {\large Rodrigo Guerreiro n\textsuperscript{o} 75739 $\cdot$ Dinis Custódio n\textsuperscript{o} 79895}
\end{center}
\vspace{2pt}
{\color{azul}\hrule height 1.6pt}
\vspace{10pt}

Os pontos 1 a 4 do enunciado correspondem a tarefas cujos resultados são entregues nos
ficheiros do ZIP (\texttt{hoos.cnf} e \texttt{sat\_solver.py}) e no output do programa.
A listagem pedida no ponto 4 é \texttt{11000} ($x_1=x_2=\text{True}$,
$x_3=x_4=x_5=\text{False}$). As questões com resposta escrita são a 5 e a 6, apresentadas
de seguida.

\section{Questão 5}

\begin{quote}\itshape
Run your algorithm on the uf20-01.cnf, uf20-03.cnf, and uf20-03.cnf instances.
How many solutions satisfy each problem instance?
\end{quote}

Usámos as três primeiras instâncias: \texttt{uf20-01}, \texttt{uf20-02} e
\texttt{uf20-03}. O número de soluções de cada instância é:

\begin{table}[h]
  \centering
  \begin{tabular}{lrrrr}
    \toprule
    \textbf{Instância} & \textbf{Variáveis} & \textbf{Cláusulas} & \textbf{N\textsuperscript{o} de soluções} & \textbf{Tempo de busca (s)} \\
    \midrule
    \texttt{uf20-01.cnf} & 20 & 91 & 8  & 3,14 \\
    \texttt{uf20-02.cnf} & 20 & 91 & 29 & 3,02 \\
    \texttt{uf20-03.cnf} & 20 & 91 & 1  & 3,02 \\
    \bottomrule
  \end{tabular}
  \caption{Número de soluções e tempo de busca medido pelo programa.}
\end{table}

Cada execução visita $2^{20} = 1\,048\,576$ atribuições e avalia 91 cláusulas por
atribuição com curto-circuito, isto é, a avaliação de uma cláusula para no primeiro
literal verdadeiro e uma atribuição é descartada na primeira cláusula falsa. O tempo
de busca é medido pelo próprio programa (ver Questão 6). Por exemplo, a instância
\texttt{uf20-03} tem uma única solução:

\begin{lstlisting}
11110111111010011101
\end{lstlisting}

\section{Questão 6}

\begin{quote}\itshape
It is impractical to do the same for a larger instance, say one contained in uf100-430.
How could you estimate the time needed by your computer program to complete the task on
such an instance? Justify your answer.
\end{quote}

O programa visita exaustivamente todas as $2^n$ atribuições, pelo que o tempo cresce
exponencialmente com o número de variáveis $n$. O tempo da fase de busca é medido com
\texttt{time.process\_time()} e reportado na linha de timing do formato DIMACS
(\texttt{satformat.pdf}, secção 2.4):

\begin{lstlisting}
t cnf <SOLUTION> <VARIABLES> <CLAUSES> <CPUSECS> <MEASURE1>
\end{lstlisting}

em que \texttt{MEASURE1} $= 2^n$ é o número de atribuições visitadas --- uma medida
algorítmica e reprodutível, independente da máquina. Exemplo real (\texttt{uf20-01}):

\begin{lstlisting}
t cnf 1 20 91 3.141 1048576
\end{lstlisting}

A estimativa para $n = 100$ faz-se extrapolando esse valor medido:

\begin{itemize}[leftmargin=1.4em,itemsep=2pt]
  \item \texttt{uf100-430} tem $n = 100$ variáveis, logo
        $2^{100} \approx 1{,}27 \times 10^{30}$ atribuições.
  \item Fator de crescimento face a $n = 20$:
        $2^{100}/2^{20} = 2^{80} \approx 1{,}21 \times 10^{24}$.
  \item Tempo medido para $n = 20$: $T(20) \approx 3{,}1$ s
        (mínimo de 3 execuções; os tempos variam com a carga da máquina).
  \item $T(100) \approx 3{,}1 \times 2^{80}$ s
        $\approx 3{,}7 \times 10^{24}$ s
        $\approx 1{,}2 \times 10^{17}$ anos.
\end{itemize}

\textbf{Justificação:}

\begin{itemize}[leftmargin=1.4em,itemsep=2pt]
  \item A complexidade é $\Theta(2^n \cdot m)$. O corte ao primeiro literal/cláusula
        falsa (curto-circuito) e o aumento do número de cláusulas
        ($91 \rightarrow 430$) são apenas fatores constantes e não alteram o
        crescimento exponencial.
  \item Mesmo com um computador $10^9$ vezes mais rápido, o tempo continuaria a ser
        cerca de $1{,}2 \times 10^{8}$ anos, pelo que a força bruta é inviável para
        $n = 100$.
  \item Para instâncias maiores seria necessário substituir a enumeração exaustiva por
        métodos que explorem a estrutura do problema, como retrocesso com propagação
        (DPLL), \textit{solvers} CDCL ou heurísticas de procura local no contexto MAXSAT.
\end{itemize}

\end{document}
